\textbf{Phần 1. Trắc nghiệm lựa chọn}
\setcounterpageref{ex}{0}

\begin{ex} %[1D1V05]
Số 7 gồm 4 và mấy?
\choice
{$2$}
{$4$}
{\True $3$}
{$1$}
\loigiai{
Ta có sơ đồ tách số: $7$ gồm $4$ và $3$ (vì $4 + 3 = 7$).
\textbf{Đáp số:} $3$.\par
\textbf{Đáp án của câu hỏi là C}
}
\end{ex}

\begin{ex} %[1D1V05]
Số 9 gồm 2 và mấy?
\choice
{$6$}
{\True $7$}
{$8$}
{$5$}
\loigiai{
Ta có sơ đồ tách số: $9$ gồm $2$ và $7$ (vì $2 + 7 = 9$).
\textbf{Đáp số:} $7$.\par
\textbf{Đáp án của câu hỏi là B}
}
\end{ex}

\begin{ex} %[1D1V05]
Số 10 gồm mấy và 6?
\choice
{$3$}
{$5$}
{$2$}
{\True $4$}
\loigiai{
Ta có sơ đồ tách số: $10$ gồm $4$ và $6$ (vì $4 + 6 = 10$).
\textbf{Đáp số:} $4$.\par
\textbf{Đáp án của câu hỏi là D}
}
\end{ex}

\begin{ex} %[1D1V05]
Gộp 5 và 3 được số nào?
\choice
{\True $8$}
{$7$}
{$9$}
{$6$}
\loigiai{
Ta có sơ đồ gộp số: Gộp $5$ và $3$ được: $5 + 3 = 8$.
\textbf{Đáp số:} $8$.\par
\textbf{Đáp án của câu hỏi là A}
}
\end{ex}

\begin{ex} %[1D1V05]
Gộp 4 và mấy để được 9?
\choice
{$4$}
{$6$}
{\True $5$}
{$3$}
\loigiai{
Để gộp 4 với một số được 9, số đó phải là: $9 - 4 = 5$.
\textbf{Đáp số:} $5$.\par
\textbf{Đáp án của câu hỏi là C}
}
\end{ex}

\begin{ex} %[1D1V05]
Có bao nhiêu cách tách số 5 thành hai số tự nhiên khác 0 (không tính thứ tự đảo ngược)?
\choice
{$1$}
{\True $2$}
{$3$}
{$4$}
\loigiai{
Các cách tách số 5 thành hai số khác 0 là:\n- $5 = 1 + 4$\n- $5 = 2 + 3$\nVậy có tất cả $2$ cách tách.
\textbf{Đáp số:} $2$.\par
\textbf{Đáp án của câu hỏi là B}
}
\end{ex}

\begin{ex} %[1D1V05]
Có bao nhiêu cách tách số 6 thành hai số tự nhiên khác 0 (không tính thứ tự đảo ngược)?
\choice
{$2$}
{$4$}
{$1$}
{\True $3$}
\loigiai{
Các cách tách số 6 thành hai số khác 0 là:\n- $6 = 1 + 5$\n- $6 = 2 + 4$\n- $6 = 3 + 3$\nVậy có tất cả $3$ cách tách.
\textbf{Đáp số:} $3$.\par
\textbf{Đáp án của câu hỏi là D}
}
\end{ex}

\begin{ex} %[1D1V05]
Số 8 được tách thành hai số giống nhau. Hỏi mỗi số đó là số mấy?
\choice
{\True $4$}
{$3$}
{$5$}
{$2$}
\loigiai{
Vì hai số giống nhau có tổng bằng 8, mà $4 + 4 = 8$ nên mỗi số đó là số $4$.
\textbf{Đáp số:} $4$.\par
\textbf{Đáp án của câu hỏi là A}
}
\end{ex}

\begin{ex} %[1D1V05]
Số 10 được tách thành hai số giống nhau. Hỏi mỗi số đó là số mấy?
\choice
{$4$}
{$6$}
{\True $5$}
{$3$}
\loigiai{
Vì hai số giống nhau có tổng bằng 10, mà $5 + 5 = 10$ nên mỗi số đó là số $5$.
\textbf{Đáp số:} $5$.\par
\textbf{Đáp án của câu hỏi là C}
}
\end{ex}

\begin{ex} %[1D1V05]
Bạn Lan có 7 cái kẹo, Lan chia vào hai túi sao cho túi thứ nhất có 3 cái kẹo. Hỏi túi thứ hai có bao nhiêu cái kẹo?
\choice
{$3$}
{\True $4$}
{$5$}
{$2$}
\loigiai{
Số kẹo ở túi thứ hai là: $7 - 3 = 4$ (cái kẹo).
\textbf{Đáp số:} $4$.\par
\textbf{Đáp án của câu hỏi là B}
}
\end{ex}

\begin{ex} %[1D1V05]
Một cây có 9 quả táo. Gió thổi làm rụng 4 quả xuống đất. Hỏi trên cây còn lại bao nhiêu quả táo?
\choice
{$4$}
{$6$}
{$3$}
{\True $5$}
\loigiai{
Tách 9 quả táo thành 4 quả rụng và số quả còn lại trên cây: $9 - 4 = 5$ (quả táo).
\textbf{Đáp số:} $5$.\par
\textbf{Đáp án của câu hỏi là D}
}
\end{ex}

\begin{ex} %[1D1V05]
Gộp 3 chiếc bút chì đỏ và 6 chiếc bút chì xanh được tất cả bao nhiêu chiếc bút chì?
\choice
{\True $9$}
{$8$}
{$10$}
{$7$}
\loigiai{
Gộp số bút chì đỏ và số bút chì xanh: $3 + 6 = 9$ (chiếc bút chì).
\textbf{Đáp số:} $9$.\par
\textbf{Đáp án của câu hỏi là A}
}
\end{ex}

\textbf{Phần 2. Trắc nghiệm đúng sai}
\setcounterpageref{ex}{0}

\begin{ex} %[1D1V05]
Xét tính đúng hoặc sai của các khẳng định toán học sau đây:
\choiceTF
{\True Kết quả của bài toán «Số 6 gồm 1 và một số. Số đó là số mấy» là $5$.}
{Kết quả của bài toán «Số 8 gồm 5 và mấy» là $2$.}
{Có 10 quả cam chia đều vào hai đĩa. mỗi đĩa có $4$ quả cam.}
{\True Một nhóm gồm 8 bạn học sinh được tách thành hai nhóm nhỏ: nhóm nam và nhóm nữ. Biết nhóm nữ có 3 bạn, hỏi nhóm nam có $5$ bạn.}
\loigiai{
\begin{itemize}
    \item Ý a) Đúng: Ta có: $6$ gồm $1$ và $5$ (vì $1 + 5 = 6$).
\textbf{Đáp số:} $5$.
    \item Ý b) Sai: Kết quả đúng là $3$, không phải $2$. Ta có: $8$ gồm $5$ và $3$ (vì $5 + 3 = 8$).
\textbf{Đáp số:} $3$.
    \item Ý c) Sai: Kết quả đúng là $5$, không phải $4$. Chia đều 10 quả cam vào 2 đĩa: mỗi đĩa có $5$ quả cam (vì $5 + 5 = 10$).
\textbf{Đáp số:} $5$.
    \item Ý d) Đúng: Số bạn học sinh nam là: $8 - 3 = 5$ (bạn).
\textbf{Đáp số:} $5$.
\end{itemize}
\begin{itemize}
    \item a) \textbf{Đúng}.
    \item b) \textbf{Sai}.
    \item c) \textbf{Sai}.
    \item d) \textbf{Đúng}.
\end{itemize}
}
\end{ex}

\begin{ex} %[1D1V05]
Xét tính đúng hoặc sai của các bài toán và phép tính sau đây:
\choiceTF
{Kết quả của bài toán «Gộp 2 que tính và 7 que tính được tất cả bao nhiêu que tính» là $8$.}
{\True Số 9 gồm hai số tự nhiên liên tiếp. số lớn hơn trong hai số đó là số $5$ .}
{\True Số 7 gồm hai số tự nhiên liên tiếp. số bé hơn trong hai số đó là số $3$ .}
{Có 10 quả cam chia đều vào hai đĩa. mỗi đĩa có $4$ quả cam.}
\loigiai{
\begin{itemize}
    \item Ý a) Sai: Kết quả đúng là $9$, không phải $8$. Gộp 2 và 7 que tính: $2 + 7 = 9$ (que tính).
\textbf{Đáp số:} $9$.
    \item Ý b) Đúng: Hai số tự nhiên liên tiếp có tổng bằng 9 là $4$ và $5$ (vì $4 + 5 = 9$ và $5 - 4 = 1$).\nSố lớn hơn trong hai số đó là số $5$.
\textbf{Đáp số:} $5$.
    \item Ý c) Đúng: Hai số tự nhiên liên tiếp có tổng bằng 7 là $3$ và $4$ (vì $3 + 4 = 7$).\nSố bé hơn trong hai số đó là số $3$.
\textbf{Đáp số:} $3$.
    \item Ý d) Sai: Kết quả đúng là $5$, không phải $4$. Chia đều 10 quả cam vào 2 đĩa: mỗi đĩa có $5$ quả cam (vì $5 + 5 = 10$).
\textbf{Đáp số:} $5$.
\end{itemize}
\begin{itemize}
    \item a) \textbf{Sai}.
    \item b) \textbf{Đúng}.
    \item c) \textbf{Đúng}.
    \item d) \textbf{Sai}.
\end{itemize}
}
\end{ex}

\textbf{Phần 3. Trả lời ngắn}
\setcounterpageref{ex}{0}

\begin{ex} %[1D1V05]
Gộp 1, 3 và 5 được số nào?
\shortans{9}
\loigiai{
\textbf{Đáp số:} $9$\par
Gộp ba số lại: $1 + 3 + 5 = 9$.
\textbf{Đáp số:} $9$.
}
\end{ex}

\begin{ex} %[1D1V05]
Số 10 gồm 3 số: 2, 5 và một số nữa. Hỏi số còn lại là số mấy?
\shortans{3}
\loigiai{
\textbf{Đáp số:} $3$\par
Tổng của hai số đã biết là: $2 + 5 = 7$.\nSố còn lại là: $10 - 7 = 3$.
\textbf{Đáp số:} $3$.
}
\end{ex}

\begin{ex} %[1D1V05]
Có 8 bông hoa cắm vào hai lọ sao cho số hoa ở hai lọ bằng nhau. Hỏi mỗi lọ có bao nhiêu bông hoa?
\shortans{4}
\loigiai{
\textbf{Đáp số:} $4$\par
Vì số hoa ở hai lọ bằng nhau và $4 + 4 = 8$ nên mỗi lọ có $4$ bông hoa.
\textbf{Đáp số:} $4$.
}
\end{ex}

\begin{ex} %[1D1V05]
Tách số 7 thành hai số có hiệu bằng 3. Hỏi số lớn hơn là số mấy?
\shortans{5}
\loigiai{
\textbf{Đáp số:} $5$\par
Hai số có tổng bằng 7 và hiệu bằng 3 là $5$ và $2$ (vì $5 + 2 = 7$ và $5 - 2 = 3$).\nSố lớn hơn là số $5$.
\textbf{Đáp số:} $5$.
}
\end{ex}

\begin{ex} %[1D1V05]
Tách số 6 thành hai số có hiệu bằng 4. Hỏi số bé hơn là số mấy?
\shortans{1}
\loigiai{
\textbf{Đáp số:} $1$\par
Hai số có tổng bằng 6 và hiệu bằng 4 là $5$ và $1$ (vì $5 + 1 = 6$ và $5 - 1 = 4$).\nSố bé hơn là số $1$.
\textbf{Đáp số:} $1$.
}
\end{ex}

\begin{ex} %[1D1V05]
Gộp 4 bạn nam và 6 bạn nữ lại thành một đội văn nghệ. Hỏi đội văn nghệ đó có tất cả bao nhiêu bạn?
\shortans{10}
\loigiai{
\textbf{Đáp số:} $10$\par
Gộp số bạn nam và nữ: $4 + 6 = 10$ (bạn).
\textbf{Đáp số:} $10$.
}
\end{ex}
